\section{Esqueleto funcional}

\subsection{Qué demuestra el esqueleto}

El esqueleto no implementa la lógica de negocio: demuestra que la base técnica distribuida
funciona. Dos servicios independientes, cada uno en su proceso y su puerto, se comunican por el
mecanismo elegido en el ADR-002: el servicio de incidencias consulta al de técnicos y devuelve la
respuesta combinada.

\begin{table}[H]
\centering
\caption{Componentes del esqueleto}
\label{tab:esqueleto}
\small
\begin{tabularx}{\linewidth}{@{}lllY@{}}
\toprule
Servicio & Puerto & Rol & Endpoint demostrativo \\
\midrule
\texttt{servicio-tecnicos} & 8081 & Proveedor & \texttt{GET /api/tecnicos}, \texttt{GET /api/tecnicos/\{id\}},
\texttt{GET /actuator/health} \\
\addlinespace
\texttt{servicio-incidencias} & 8082 & Consumidor & \texttt{GET /api/incidencias/\{id\}} (resuelve el
técnico llamando al puerto 8081) \\
\bottomrule
\end{tabularx}
\end{table}

\subsection{Ejecución}

Las instrucciones completas están en el \texttt{README.md} de la raíz; en resumen, desde una
terminal por servicio:

\begin{verbatim}
cd src/servicio_tecnicos    && python app.py    # http://localhost:8081
cd src/servicio_incidencias && python app.py    # http://localhost:8082
curl -i http://localhost:8082/api/incidencias/1
\end{verbatim}

\pendiente{Agregar aquí la salida real de la verificación, con fecha de la ejecución.}

\subsection{Evidencia de ejecución}

\begin{figure}[H]
\centering
\figurapendiente{Captura 1: ambos servicios en marcha, con usuario, máquina, fecha y hora visibles}
\caption{Esqueleto en ejecución: servicios activos y llamada entre ellos.}
\label{fig:captura1}
\end{figure}

\begin{figure}[H]
\centering
\figurapendiente{Captura 2: respuesta combinada del consumidor con los datos del proveedor}
\caption{Respuesta del servicio de incidencias con datos resueltos del servicio de técnicos.}
\label{fig:captura2}
\end{figure}

\begin{figure}[H]
\centering
\figurapendiente{Captura 3: degradación controlada con el proveedor detenido (503, no 500)}
\caption{Comportamiento ante la caída del servicio proveedor: degradación controlada.}
\label{fig:captura3}
\end{figure}

\section{Aportes del equipo}

Los enlaces apuntan a las confirmaciones y a los issues reales; la tabla se actualiza el último día
con los números definitivos.

\begin{table}[H]
\centering
\caption{Aportes por integrante}
\label{tab:aportes}
\small
\begin{tabularx}{\linewidth}{@{}lYl@{}}
\toprule
Integrante & Aporte principal & Enlaces \\
\midrule
Angel (\texttt{@iMouraad}) & Arquitectura, ADR, C4, CAP/PACELC, documento & \pendiente{commits e issues} \\
Angello (\texttt{@Amendozab5}) & Esqueleto, comunicación, ejecución, capturas & \pendiente{commits e issues} \\
César (\texttt{@CesarMacias007}) & Requisitos, riesgos, pruebas, coherencia & \pendiente{commits e issues} \\
\bottomrule
\end{tabularx}
\end{table}

\section{Declaración de uso de inteligencia artificial}

\pendiente{Completar con la verdad, herramienta por herramienta. Formato sugerido:}

\begin{table}[H]
\centering
\caption{Uso de herramientas de inteligencia artificial}
\small
\begin{tabularx}{\linewidth}{@{}llYY@{}}
\toprule
Herramienta y versión & Para qué se usó & Qué se verificó & Cómo se verificó \\
\midrule
\pendiente{} & \pendiente{} & \pendiente{} & \pendiente{} \\
\bottomrule
\end{tabularx}
\end{table}

% ---------------------------------------------------------------------
% Si no se usó IA, se escribe una sola frase indicándolo y se elimina
% la tabla. Si se usó, se declara también lo usado en los documentos
% previos (FORM-AC-01 declaró ChatGPT y Claude). No declararlo es un
% problema de integridad, no de estilo.
% ---------------------------------------------------------------------
